\subsection{晶体对 X 射线的衍射}
\src{上册 §4.4，书内 p.244—253 / PDF p.263—272}

多晶衍射实验最终只给出两列数：每个峰的衍射角 $2\theta$ 与它的强度 $I$。本节余下内容的全部
推论都建立在一个分工上——\textbf{峰位}只由晶胞的平移周期决定，因而给出晶面间距与点阵参数；
\textbf{峰强}由晶胞内原子的种类、坐标、占位与热振动决定；\textbf{峰的有无}（系统消光）由点阵
类型与微观对称元素决定。三条线索彼此独立又互相印证，是物相鉴定、相定量、指标化乃至结构精修
的公共基础。

\begin{probout}
\textbf{测什么}：一套 $(2\theta,\,I)$ 峰表，以及由此换算的晶面间距 $d$ 与相对强度 $I/I_1$。\\
\textbf{能推断什么}：$d$ 序列 $\rightarrow$ 晶胞参数与点阵类型（配合系统消光）；强度分布
$\rightarrow$ 晶胞内原子的位置、种类、占位度与热振动；整套图谱 $\rightarrow$ 物相种类，多相样品
中各相相对含量。\\
\textbf{不能推断什么}：衍射实验中相位信息丢失，实测只给出 $|F_{hkl}|$，因此粉末强度本身不足以
唯一确定原子坐标；轻原子（$Z$ 小）与低含量相（$\lesssim 3\%$）在强度上贡献很小，属于方法本身的
边界而非操作失误。
\end{probout}

\subsubsection{衍射峰的位置与晶面间距}

\begin{principle}
衍射出现的必要条件（布拉格方程）
\begin{equation}
2\,d_{hkl}\sin\theta=\lambda ,
\label{eq:ch04b-bragg}
\end{equation}
它把可观测量 $\theta$ 与结构量 $d_{hkl}$ 联系起来。晶面间距由晶格常数 $a,b,c$ 及三个夹角
$\alpha,\beta,\gamma$（或其倒易点阵常数 $a^{*},b^{*},c^{*}$ 与 $\alpha^{*},\beta^{*},\gamma^{*}$）
和衍射指标 $h,k,l$ 共同决定。以倒易格子向量 $\vec{H}=h\vec{a}^{*}+k\vec{b}^{*}+l\vec{c}^{*}$
为工具，可得
 
\begin{equation}
Q\equiv\frac{1}{d_{hkl}^{2}}=\left|\vec{H}\right|^{2},
\label{eq:ch04b-Q}
\end{equation}
即 $Q$ 是倒易向量长度的平方。$Q$ 的引入使“峰位 $\rightarrow$ 结构”的关系在一个晶系内变成线性
关系：立方、四方、六方、正交晶系中 $Q$ 都是 $h,k,l$ 平方项之和。此处已对照原页面图像
（书内 p.244—245 / PDF p.263—264）：原书\textbf{并未印出通式} $Q\equiv1/d^{2}=|\vec{H}|^{2}$，
而是逐晶系直接给出 $1/d^{2}$ 的表达式——p.244 给立方、四方、正交、六方四式，p.245 给单斜、
三斜两式，并另行给出 R 格子取素晶格（$a=b=c$、$\alpha=\beta=\gamma$）时的简化式
$1/d^{2}=(h^{2}+k^{2}+l^{2})a^{*2}+2(hk+kl+hl)a^{*2}\cos\alpha^{*}$。原书只说明「利用倒易格子向量
$\vec{H}=h\vec{a}^{*}+k\vec{b}^{*}+l\vec{c}^{*}$ 及其与晶面间距之间的关系，可以求出……数学关系」。
故以 $\vec{H}$ 长度平方定义 $Q$、把各晶系写成统一的 $Q$ 记号，是本书为把「峰位 $\rightarrow$ 结构」
写成线性关系而引入的组织方式，不是原书原符号；下表各式的数学内容与原书逐式等价，只是原书
单斜、三斜、菱方形用倒易点阵常数表达，本表改写为直接晶格常数的等价式。
\end{principle}

\begin{table}[tbp]
  \centering
  \caption{各晶系的 $Q=1/d^{2}$ 表达式的标准形式，以及需要确定的独立晶胞参数个数。单斜式取
    $b$ 轴唯一（$\beta$ 为 $a$ 与 $c$ 的夹角）。\src{据 S1（上册）§4.4.1，书内 p.244—245 / PDF p.263—264 归纳。}}
  \label{tab:ch04b-Qtable}
  \begin{tabular}{@{}llc@{}}
    \toprule
    晶系 & $Q=1/d^{2}$ & 独立参数 \\
    \midrule
    立方 & $\dfrac{h^{2}+k^{2}+l^{2}}{a^{2}}$ & 1 \\
    四方 & $\dfrac{h^{2}+k^{2}}{a^{2}}+\dfrac{l^{2}}{c^{2}}$ & 2 \\
    六方（三方取六方轴） & $\dfrac{4(h^{2}+hk+k^{2})}{3a^{2}}+\dfrac{l^{2}}{c^{2}}$ & 2 \\
    菱方（取菱方轴） & $\dfrac{(h^{2}+k^{2}+l^{2})\sin^{2}\alpha+2(hk+kl+lh)(\cos^{2}\alpha-\cos\alpha)}
      {a^{2}(1-3\cos^{2}\alpha+2\cos^{3}\alpha)}$ & 2 \\
    正交 & $\dfrac{h^{2}}{a^{2}}+\dfrac{k^{2}}{b^{2}}+\dfrac{l^{2}}{c^{2}}$ & 3 \\
    单斜 & $\dfrac{1}{\sin^{2}\beta}\left(\dfrac{h^{2}}{a^{2}}+\dfrac{k^{2}\sin^{2}\beta}{b^{2}}
      +\dfrac{l^{2}}{c^{2}}-\dfrac{2hl\cos\beta}{ac}\right)$ & 4 \\
    \bottomrule
  \end{tabular}
\end{table}

三斜晶系（6 个独立参数）不适合写成一个紧凑分式，通常用倒易度量张量分量的形式：
\begin{equation}
Q=\frac{1}{V^{2}}\left(S_{11}h^{2}+S_{22}k^{2}+S_{33}l^{2}+2S_{12}hk+2S_{13}hl+2S_{23}kl\right),
\label{eq:ch04b-triclinic}
\end{equation}
其中 $S_{11}=b^{2}c^{2}\sin^{2}\alpha$，$S_{22}=a^{2}c^{2}\sin^{2}\beta$，$S_{33}=a^{2}b^{2}\sin^{2}\gamma$，
$S_{12}=abc^{2}(\cos\alpha\cos\beta-\cos\gamma)$，$S_{13}=ab^{2}c(\cos\alpha\cos\gamma-\cos\beta)$，
$S_{23}=a^{2}bc(\cos\beta\cos\gamma-\cos\alpha)$，
$V=abc\sqrt{1-\cos^{2}\alpha-\cos^{2}\beta-\cos^{2}\gamma+2\cos\alpha\cos\beta\cos\gamma}$。
此处已对照原页面图像（书内 p.245 / PDF p.264）：原书三斜式印为
$1/d^{2}=h^{2}a^{*2}+k^{2}b^{*2}+l^{2}c^{*2}+2hka^{*}b^{*}\cos\gamma^{*}+2klb^{*}c^{*}\cos\alpha^{*}
+2hla^{*}c^{*}\cos\beta^{*}$，单斜式印为
$1/d^{2}=h^{2}a^{*2}+k^{2}b^{*2}+l^{2}c^{*2}+2hla^{*}c^{*}\cos\beta^{*}$，均用\textbf{倒易点阵常数}
表达。式~\equref{eq:ch04b-triclinic} 的 $S_{ij}$ 记号是本项目为可读性引入的约定（等价于倒易度量
张量分量写法），与原书符号不同但数学内容一致。

\begin{measure}
\textbf{峰位怎么读}。$2\theta$ 的取值方式直接决定 $d$ 的精度：峰顶法读数受统计涨落影响最大，
重心法和全谱拟合（剖面对峰形建模后取峰位）精度最高，弱峰尤其要用拟合而不是目测峰顶。
\textbf{必须做的三项校正}：（i）剥离 $K\alpha_{2}$（$K\alpha_{1}/K\alpha_{2}$ 双线使峰不对称，
未剥离时 $d$ 系统偏大）；（ii）校正仪器零点与样品位移（样品表面偏离测角仪圆心，$d$ 的系统误差
近似正比于 $\cos\theta$，低角度最敏感）；（iii）用内标物（如 \ce{SiO2} 石英、\ce{CaF2}、Si）
把实测 $d$ 归算到标准值上。波长取 $K\alpha_{1}$ 单线时 Cu 靶为 $1.5406$ Å，取 $K\alpha$ 加权
平均值时为 $1.5418$ Å；两者混用会引入约 $0.05\%$ 的系统偏差。
\end{measure}

\begin{dataflow}
$2\theta$（已剥离 $K\alpha_{2}$、已做位移与内标校正）$\rightarrow$ 由式~\equref{eq:ch04b-bragg} 得
$d_{hkl}$ $\rightarrow$ 由式~\equref{eq:ch04b-Q} 得 $Q$ $\rightarrow$ 对每个可能的指标 $(hkl)$
写出一个线性方程 $\rightarrow$ 最小二乘解出晶胞参数 $\rightarrow$ 高角度线权重更大（$\Delta d/d$
随 $\theta$ 增大而减小），故晶胞参数精修应优先采用高角线并作外推。
\end{dataflow}

\begin{pitfall}
把峰位误差当成结构信息的三种典型情形：\\
（1）\textbf{样品位移}造成所有峰的 $d$ 同向偏移，看起来像“晶胞整体胀大”。判别方法：加内标后
偏移是否消失；或看 $\Delta d/d$ 是否随 $\cos\theta$ 规律变化而非随机。\\
（2）\textbf{未剥离 $K\alpha_{2}$}使每个峰出现“肩”，重心法读数介于两条线之间，$d$ 偏大且偏移
随角度变化，容易被误读为固溶导致的晶胞参数变化。\\
（3）\textbf{峰位随温度漂移}（热膨胀）。同一批样品在不同温度或不同时间测量时，$d$ 的小幅变化
要先排除温度因素，再讨论成分或缺陷。
\end{pitfall}

\begin{crosscheck}
峰位得到的晶胞参数必须与指标化结果（见 §4.6）互校：若一套指标能给出一致的晶胞参数且能指标化
全部衍射线，则峰位数据自洽；若不能，往往是指标化选错或存在未识别的第二相。进一步可与 PDF
卡片值、单晶结构数据、以及化学分析给出的元素比互校。
\end{crosscheck}

\subsubsection{多晶衍射峰的强度}

\begin{principle}
多晶（粉末）样品的衍射强度是“单晶某一 $hkl$ 的散射能力”乘上一系列几何与物理修正后的结果。
本书采用的简化形式为
\begin{equation}
I_{hkl}=K\cdot\left|F_{hkl}\right|^{2}\cdot P\cdot L\cdot j\cdot D\cdot\Delta V,
\qquad
K=I_{0}\left(\frac{e^{4}}{m^{2}c^{4}}\right)\frac{\lambda^{3}}{8\pi R^{2}V_{c}^{2}},
\label{eq:ch04b-intensity}
\end{equation}
式中 $I_{0}$ 为入射线强度，$e$ 为电子电荷，$m$ 为电子质量，$c$ 为光速，$\lambda$ 为波长，$R$ 为
样品中心至衍射图（探测器）的距离即衍射仪半径，$V_{c}$ 为晶胞体积；$\theta$ 为布拉格角，
$P$ 为偏振因子，$L$ 为洛伦兹因子，$F_{hkl}$ 为结构因子，$j$ 为多重度因子，$D$ 为温度因子
（Debye 因子），$\Delta V$ 为参与衍射的样品体积。此处常数项已逐项对照原页面图像核对：
原书印
$I_{hkl}=I_{0}(e^{4}/m^{2}c^{4})\cdot\lambda^{3}/(8\pi R^{2}V_{c}^{2})\cdot
[(1+\cos^{2}2\theta)/2]\cdot[1/(\sin^{2}\theta\cos\theta)]\cdot|F_{hkl}|^{2}\cdot D\cdot\Delta V\cdot j$，
并写明上述参数「可以合并为常数项 $K'$」；本节把 $3$ 个角度因子分列，故写作 $P$、$L$、$j$ 的
乘积形式（原书 $K'$ 与本式 $K$ 同义）。\textbf{原书参数表中没有独立的吸收因子 $A$}：吸收效应
被折进体积因子 $\Delta V$——原书印「当考虑吸收效应时，样品参与衍射的有效体积将小于 $\Delta V$
……而平行板样品吸收效应与角度无关，$\Delta V$ 可以 $A_{0}/2\mu$ 代之」。
\end{principle}

\begin{equation}
P=\frac{1+\cos^{2}2\theta}{2},\qquad
L=\frac{1}{2\sin^{2}\theta\cos\theta},\qquad
PL=\frac{1+\cos^{2}2\theta}{4\sin^{2}\theta\cos\theta}.
\label{eq:ch04b-PL}
\end{equation}
三式已逐项对照原页面图像核对：原书印
$P=(1+\cos^{2}2\theta)/2$、$L=1/(2\sin^{2}\theta\cos\theta)$、
$PL=(1+\cos^{2}2\theta)/(4\sin^{2}\theta\cos\theta)$，与本式逐字一致，分母确为
$4\sin^{2}\theta\cos\theta$。单晶／de Wolff 惯例中常把洛伦兹因子写作
$L=1/(\sin^{2}\theta\cos\theta)$，与原书相差一个可并入常数项的因子 $2$；本节取原书的刻度。

\begin{equation}
D=\exp\!\left(-\frac{B\sin^{2}\theta}{\lambda^{2}}\right),\qquad B=8\pi^{2}U,\qquad U=\langle u^{2}\rangle ,
\label{eq:ch04b-DW}
\end{equation}
本式已逐项对照原页面图像核对：原书即印 $D=\exp(-B\sin^{2}\theta/\lambda^{2})$，并定义
$B=8\pi^{2}U$、$U=\langle u^{2}\rangle$，「$\langle u^{2}\rangle$ 是位移振幅的均方值」——
\textbf{指数中的系数是 $B$，不是 $2B$}。原书另给量级参考：无机晶体的温度因子 $B<1$\,\AA$^{2}$，
而含 C、H、O、N 等轻原子的有机物 $B$ 可达 $2\sim6$\,\AA$^{2}$；原书未区分
$\langle u^{2}\rangle$ 取沿衍射矢量方向的投影还是三维均方值（两者相差因子 $3$），本节沿用
原书的笼统写法。
只有各原子的 $B$ 相同时，才能像本式这样把温度因子提到 $\left|F\right|^{2}$ 之外作统一乘子；
若各原子 $B_{j}$ 不同，必须逐个原子写在 $F_{hkl}$ 内部再平方，不能共用一个 $D$。
\textbf{一处口径差异记录在案}：若按「第 $j$ 个原子的散射振幅乘
$\exp[-B_{j}\sin^{2}\theta/\lambda^{2}]$、强度再取 $\left|F_{hkl}\right|^{2}$」的分解，强度层面
会得到 $\exp(-2B\sin^{2}\theta/\lambda^{2})$，与原书所印相差因子 $2$；这是原书把温度因子直接
当作强度修正因子的口径，本节\textbf{照录原书印法}而不改动公式，仅在此说明差异。

\begin{keycond}
\textbf{为什么每一项修正都不能省}：\\
$\left|F_{hkl}\right|^{2}$ 是唯一携带结构信息的项——它决定哪些峰强、哪些峰弱，是原子坐标、
占位度与元素种类的直接体现。\\
$j$（多重度因子）把“$d$ 相同、指标不同”的衍射峰合并计数。多晶样品中这些峰的衍射角完全重合，
强度叠加在一起，不乘 $j$ 就无法与单晶意义上的 $|F|^{2}$ 比较。\\
$L$（洛伦兹因子）修正衍射几何：粉末中满足布拉格条件的晶粒取向只占球面上一条环带，环带长度
与接受狭缝在 $2\theta$ 上的有效覆盖都随角度变化，参与衍射的晶粒数目因此随 $\theta$ 改变。\\
$P$（偏振因子）修正散射本身：电子散射 X 射线后射线被部分偏振，散射强度按
$(1+\cos^{2}2\theta)/2$ 衰减，$2\theta=90°$ 处损失最大（降至 $1/2$）。\\
$D$（温度因子）修正热振动：原子并不静止在坐标参数所指的位置，热振动增大原子散射波之间的
位相差，强度按指数衰减，\textbf{高角度衰减远大于低角度}。忽略它会系统性高估高角峰强度，
使由强度反推的占位度与坐标发生偏移。\\
$A$（吸收因子）修正吸收：\textbf{Bragg--Brentano 平板样品}在“无限厚”条件下
$A=1/(2\mu_{l})$ 与 $\theta$ 无关（$\mu_{l}$ 为线吸收系数），这正是衍射仪适合做相定量的根本
原因；而棒状（Debye--Scherrer）样品与有限厚平板样品 $A(\theta)=1-\exp(-2\mu_{l}t/\sin\theta)$
随 $\theta$ 变化，低角吸收强、高角吸收弱，必须单独处理。\\
$\Delta V$ 修正“有多少样品在贡献强度”。平板无限厚样品中有效体积
$\Delta V=A_{0}/(2\mu_{l})$（$A_{0}$ 为入射线垂直截面积），故强度具有\textbf{容量性质}：
被测体积越大信号越强，这一性质是相定量公式的直接来源。
\end{keycond}

\begin{table}[tbp]
  \centering
  \caption{常见晶系（取该晶系最高对称点群）的多重度因子 $j$。括号内为衍射指标类型。对称性
    降低时表中数值按点群阶次成比例减小或分裂。\src{据 S1（上册）表~4-7「对称性、衍射指标与多重度因子」，书内 p.247 / PDF p.266 整理。}}
  \label{tab:ch04b-multiplicity}
  \begin{tabular}{@{}lllll@{}}
    \toprule
    晶系（点群） & 一般 $(hkl)$ & 一个指标为零 & 两个指标相等 & 轴面 $(h00)$／$(00l)$ \\
    \midrule
    立方（$m\bar{3}m$） & 48 & 24 $(hk0)$ & 24 $(hhl)$、12 $(hh0)$ & 6、$(hhh)$ 8 \\
    六方／三方（$6/mmm$） & 24 & 12 $(h0l)$、12 $(hk0)$ & 12 $(hhl)$、6 $(hh0)$ & 6、2 \\
    四方（$4/mmm$） & 16 & 8 $(h0l)$、8 $(hk0)$ & 8 $(hhl)$、4 $(hh0)$ & 4、2 \\
    正交（$mmm$） & 8 & 4 $(h0l)$、4 $(hk0)$、4 $(0kl)$ & --- & 2、2、2 \\
    单斜（$2/m$） & 4 & 2 $(h0l)$、2 $(hk0)$ & --- & --- \\
    三斜（$\bar{1}$） & 2 & 2 & 2 & 2 \\
    \bottomrule
  \end{tabular}
\end{table}

表~\tabref{tab:ch04b-multiplicity} 的用法是：多晶谱上某峰的实测强度，等于该 $Q$ 值下所有
“$h,k,l$ 置换等价”晶面的 $\left|F\right|^{2}$ 之和乘以其余因子。例如正交晶系的 $hkl$
（$h,k,l$ 互不相等且均不为零）有 8 个等价晶面，若它们 $\left|F\right|^{2}$ 相等，则实测强度
是单个 $\left|F\right|^{2}$ 的 8 倍。
本表已逐格对照原页面图像核对：原书表 4-7 列出立方 $O_h\,O\,T_d$（48、24 $hhl$、24 $hk0$、
12 $hh0$、8 $hhh$、6 $h00$）、六方／三方 $D_{6h}\,D_6\,D_{6v}\,D_{3h}$（24、12、12、12、6、
6、2）、四方 $D_{4h}\,D_4\,C_{4v}\,D_{2d}$（16、8、8、8、4、4、2）、正交
$D_{2h}\,D_2\,C_{2v}$（8、4、4、4、2、2、2）、单斜 $C_{2h}\,C_2\,C_s$（4、2、2）与三斜
$C_i\,C_1$（所有类型 $j=2$）各行，本表所取值与之一致；原书同一晶系内再按点群分列，低对称行
的数值以「$\times 2$」形式分裂（如 $T_h\,T$ 行为 $24\times2$、$12\times2$），本表按该晶系
最高对称点群压缩为一行，并在括号中保留衍射指标类型。原书把 $hk0$ 单列，本表已并入
「一个指标为零」列。\src{S1 表~4-7，书内 p.247 / PDF p.266}

\begin{example}[读谱实例：择优取向造成的强度异常]
\begin{obs}
一组示意数据（模型构造，非实测）：某六方层状氧化物用两种制样各测一次——(a) 常规正压片；
(b) 侧装法（压制方向与样品面平行）装样，其余实验条件完全相同。测得积分强度比
$I_{002}/I_{100}$：(a) 为 $12.5$，(b) 为 $1.8$。两种制样下全部峰的 $2\theta$ 峰位一致，
偏差小于 $0.02°$。
\end{obs}
\begin{infer}
峰位一致说明晶胞参数与相组成没有变化；而 $I_{002}/I_{100}$ 随制样方式改变近 7 倍，说明强度
差异来自晶粒取向的统计分布（择优取向），不是原子占位或结构的改变。若把 (a) 的异常强度当成
“结构变化”去解释——例如认为层间插入了重原子或阳离子占位发生了迁移——就会得到错误结论。
\end{infer}
正确的处理顺序是：先在改变制样方式（侧装、轻压、加冲淡剂、减小粒度）后重测，看强度比是否
随之改变；只有\textbf{在取向被消除后仍然存在}的强度异常，才值得用结构模型去解释。
\end{example}

\begin{figure}[tbp]
  \centering
  \begin{tikzpicture}
    \tikzset{declare function={
        PLf(\x) = (1+cos(\x)^2)/(4*sin(\x/2)^2*cos(\x/2));
        DWf(\x,\B) = exp(-2*\B*sin(\x/2)^2/2.3734);
      }}
    \begin{axis}[
        width=\onecolwidth, height=64mm,
        xlabel={$2\theta$ / (°)},
        ylabel={相对因子（对数纵轴）},
        xmin=10, xmax=160, ymin=0.01, ymax=1.7,
        ymode=log,
        xtick={10,30,50,70,90,110,130,150},
        ytick={0.02,0.05,0.1,0.2,0.5,1},
        yticklabels={0.02,0.05,0.1,0.2,0.5,1},
        grid=both, major grid style={gray!35}, minor grid style={gray!12},
        tick label style={font={\fontsize{8}{9.6}\selectfont}},
        label style={font={\fontsize{8}{9.6}\selectfont}},
        legend style={font={\fontsize{8}{9.6}\selectfont}, draw=none,
                      fill=white, fill opacity=0.9, text opacity=1,
                      legend cell align=left,
                      at={(0.985,0.975)}, anchor=north east},
        samples=201, domain=10:160,
      ]
      \addplot[very thick, blue] {PLf(x)/PLf(10)};
      \addlegendentry{$PL$ 因子（$2\theta=10°$ 处归一为 1）}
      \addplot[very thick, red, dashed] {DWf(x,1)};
      \addlegendentry{温度因子 $D$，$B=1$ Å$^2$}
      \addplot[very thick, orange!80!black, dotted] {DWf(x,3)};
      \addlegendentry{温度因子 $D$，$B=3$ Å$^2$}
      \addplot[very thick, black!65] {PLf(x)/PLf(10)*DWf(x,1)};
      \addlegendentry{$PL\cdot D$（$B=1$ Å$^2$）}
    \end{axis}
  \end{tikzpicture}
  \caption{两个必须做的角度修正因子随 $2\theta$ 的变化（模型计算，按式~\equref{eq:ch04b-PL}
    与式~\equref{eq:ch04b-DW}，Cu $K\alpha$，$\lambda=1.5406$ Å）。\textbf{看什么}：$PL$ 因子
    在低角极大、随角度迅速跌落（$2\theta=10°$ 处约 65，$2\theta=160°$ 处约 2.8，正好落在原书
    PL 因子图纵轴 0—80 的量程内，可作为公式重建的旁证）；温度因子在低角几乎为 1，在高角明显
    衰减，$B$ 越大衰减越强（$2\theta=160°$ 时 $B=1$ Å$^2$ 降到约 $0.44$，$B=3$ Å$^2$ 降到约
    $0.086$）。\textbf{结论}：把两条曲线相乘后，强度的角度依赖比任何单项都剧烈，因此低角与高角
    的实测强度不能直接互比，必须先做 $PL$ 与 $D$ 修正。\textbf{局限}：纵轴为对数，曲线是相对值；
    $B$ 的取值是示意（无机晶体常 $B<1$ Å$^2$，含 C、H、O、N 的有机物可到 $2$—$5$ Å$^2$）。
    \figsrc{模型计算图，非原书插图；$PL$ 的形式与原书第 246 页 PL 因子图一致。}}
  \label{fig:ch04b-pld}
\end{figure}

\begin{figure}[tbp]
  \centering
  \begin{tikzpicture}
    \begin{axis}[
        width=\onecolwidth, height=64mm,
        ybar, bar width=6mm,
        xmin=0.45, xmax=8.55, ymin=0, ymax=3.6,
        xtick={1,2,3,4,5,6,7,8},
        xticklabels={002,004,100,101,102,103,110,112},
        xlabel={衍射指标 $hkl$（六方层状相模型）},
        ylabel={$I_{\text{取向}}/I_{\text{随机}}$},
        ytick={0,0.5,1.0,1.5,2.0,2.5,3.0},
        grid=major, major grid style={dashed, gray!45},
        tick label style={font={\fontsize{8}{9.6}\selectfont}},
        label style={font={\fontsize{8}{9.6}\selectfont}},
      ]
      \addplot[fill=blue!45, draw=blue!70!black] coordinates {
        (1,2.37) (2,2.37) (3,0.65) (4,1.21) (5,0.82) (6,1.71) (7,0.65) (8,0.72)
      };
      \addplot[black, thick, dashed, sharp plot, update limits=false]
        coordinates {(0.45,1) (8.55,1)};
      \node[align=left, text width=62mm, anchor=north west,
            font={\fontsize{8}{9.6}\selectfont}]
        at (axis cs:0.55,3.55)
        {March--Dollase 模型，$r=0.75$：\\ $00l$ 反射增强、$hk0$ 反射减弱，\\ 虚线为随机取向水平 $1$；
         比值偏离 $1$ \\ 完全来自制样取向，与结构无关。};
    \end{axis}
  \end{tikzpicture}
  \caption{择优取向对相对强度的影响（模型计算）。\textbf{看什么}：同一物相、同一晶胞，只改变
    晶粒取向的统计分布，相对强度就可以整体偏离随机取向的水平（虚线）达数倍，且偏离方向与
    取向方向系统的关联——层状晶粒平躺在样品面上时，反射面平行于样品面的 $00l$ 被增强，
    垂直的 $hk0$ 被减弱。\textbf{条件}：$r$ 为 March--Dollase 取向参数（$r=1$ 为随机取向），
    各峰与取向轴的夹角取 $(002,004)=0°$、$(103)=25°$、$(101)=40°$、$(102)=60°$、
    $(112)=70°$、$(100,110)=90°$；柱高为取向样品与随机取向样品的强度比。
    \textbf{局限}：这是单一取向轴的理想模型，实际样品常同时存在两种取向分量，需要用
    Rietveld 精修中的取向模型拟合，而不能靠单峰强度比判断。\figsrc{模型计算图，非原书插图。}}
  \label{fig:ch04b-texture}
\end{figure}

\begin{dataflow}
\textbf{由强度到结构信息}：实测积分强度 $\rightarrow$ 扣背底、剥离 $K\alpha_{2}$、分离重叠峰
$\rightarrow$ 除以 $PL\cdot j\cdot D\cdot A$ $\rightarrow$ 得到各峰的 $\left|F_{hkl}\right|^{2}$
$\rightarrow$ 与由结构模型（原子坐标、占位度、$B$）计算的 $\left|F_{hkl}\right|^{2}$ 比较
$\rightarrow$ 修正模型使全部峰的残差最小。注意这条链路给出的是 $\left|F\right|^{2}$ 的相对值，
只有经过标定（内标或已知结构参照）才能得到绝对尺度上的占位度。
\end{dataflow}

\begin{pitfall}
\textbf{消光是另一类使实测强度低于计算值的因素}。原消光（理想完整晶体或较大镶嵌块内部入射线
与衍射线相互干涉）与次消光（镶嵌块之间的屏蔽）都使实测强度偏低，其中原消光通常影响更大，
在晶体尺寸大、完整性好、衍射强度大且衍射角小时最严重。多晶样品晶粒小，消光比单晶法弱得多，
但\textbf{低角度强线}的实测值通常仍偏小——这正是用低角强线做定量时结果偏低的来源之一。
判别线索：计算值与实测值的偏差随强度增大、角度减小而系统增大，而不是随机分布。
\end{pitfall}

\subsubsection{原子散射因子与晶体的结构因子}

原子对 X 射线的散射源于电子散射，重元素的散射能力远高于轻原子。原子中所有电子散射的
综合效应定义为原子散射因子
\begin{equation}
f=\frac{E_{a}}{E_{e}},
\label{eq:ch04b-f}
\end{equation}
$E_{a}$ 为一个原子的散射振幅，$E_{e}$ 为一个自由电子的散射振幅。若原子的全部电子集中于一点，
则 $f$ 等于电子数 $Z$；但电子分布在核周围形成电子云，不同轨道上的电子散射波之间存在位相差，
因此一般有 $f<Z$，且随散射角增大而减小（$\theta=0$ 时无位相差，$f=Z$）。定量表达为
\begin{equation}
f\!\left(\frac{\sin\theta}{\lambda}\right)=\int_{0}^{\infty}4\pi r^{2}\rho(r)\,
\frac{\sin\!\left(4\pi r\sin\theta/\lambda\right)}{4\pi r\sin\theta/\lambda}\,dr ,
\label{eq:ch04b-fint}
\end{equation}
式中 $r$ 为电子离核距离，$\rho(r)$ 为电子云密度径向分布，$\lambda$ 为 X 射线波长。
此处已对照原页面图像（书内 p.248 / PDF p.267）：原书印为
$f=\dfrac{4\pi}{e}\displaystyle\int_{0}^{\infty}r^{2}\rho(r)\,\dfrac{\sin kr}{kr}\,dr$，
$k=\dfrac{4\pi\sin\theta}{\lambda}$，即 $4\pi r^{2}$ 权重与 $\sin x/x$ 因子、以及
$k$ 与 $\sin\theta/\lambda$ 的换算关系与原书一致。差别只在归一化前置因子：原书多一个
$1/e$（与其把 $\rho(r)$ 称作「电子云密度」并列的写法），本式按标准形式（$\rho$ 取电子数
密度，$\theta=0$ 时 $f=Z$）给出，与原书正文「当散射角 $\theta=0$……$f$ 等于电子数目 $Z$」
一致。原书同页另注：$\theta=0$ 各位相差为 $0$；$\theta\neq0$ 总有相互「抵消」，故 $f<Z$。

实际计算中不用积分式，而是用拟合多项式：以 4 组高斯项加一个常数共 9 个参数表达 $f$ 随
$\sin\theta/\lambda$ 的变化，
\begin{equation}
f\!\left(\frac{\sin\theta}{\lambda}\right)=\sum_{i=1}^{4}a_{i}\exp\!\left[-b_{i}
\left(\frac{\sin\theta}{\lambda}\right)^{2}\right]+c ,
\label{eq:ch04b-fpoly}
\end{equation}
各项系数可从《国际晶体学表》C 卷表 6.1.1.4（572—574 页）查到（已对照原页面图像，书内
p.248 / PDF p.267；原书另注该多项式由 HFS（Hartree-Fock-Slater，自洽场）法与 TFD
（Thomas-Fermi-Dirac 统计）法理论计算并结合实验数据得出，$9$ 个参数与原子或离子的电子云
分布有关；同页图 4-15 给出核外电子数均为 $10$ 的原子或离子——Si$^{4+}$、Al$^{3+}$、
Mg$^{2+}$、Na$^{+}$、Ne、F$^{-}$ 及 "Point atom"——的 $f$ 随 $\sin\theta/\lambda$ 的变化）。

由 $f$ 出发给出三条读数规律：原子散射因子随衍射角增大而减小，这是高角衍射峰减弱的重要因素；
核外电子数相等的原子实，核电荷越大、核对电子吸引越强、电子云越接近点状，$f$ 越大；同一原子
的最外层电子主要贡献低角度散射，内层电子对所有角度都有较大贡献。

晶胞的结构因子定义为晶胞内全部原子散射振幅的向量叠加：
\begin{equation}
F_{hkl}=\sum_{j=1}^{n}f_{j}\exp\!\left[2\pi i\left(hx_{j}+ky_{j}+lz_{j}\right)\right],
\label{eq:ch04b-F}
\end{equation}
$n$ 为晶胞中的原子数，$(x_{j},y_{j},z_{j})$ 为第 $j$ 个原子的分数坐标，$f_{j}$ 为其原子散射
因子。$F_{hkl}$ 既表示振幅也表示相角：写作 $F=\left|F\right|\exp(i\alpha)$，$\alpha$ 为
各原子散射波的合成相角。当晶胞为素格子且只有一个原子位于原点时 $F_{hkl}=f_{A}$，相角为零；
当格子上有两个原子且第二个位于 $(x,y,z)$ 时，二者光程差与位相差为
$\Delta=(hx+ky+lz)\lambda$、$\varphi=2\pi(hx+ky+lz)$，于是
$F_{hkl}=f_{A}+f_{B}\exp\!\left[2\pi i(hx+ky+lz)\right]$，即两个散射因子的向量加和。

在反常散射不严重时，$\left|F_{hkl}\right|=\left|F_{\bar{h}\bar{k}\bar{l}}\right|$，这就是
Friedel 定律；其结果是衍射空间的强度分布具有中心对称，$230$ 个空间群只表现出 $11$ 种
Laue 群的对称性。结构因子与晶胞内电子云密度 $\rho(x,y,z)$ 之间通过傅里叶变换联系：
\begin{equation}
F_{hkl}=\int_{V_{c}}\rho(x,y,z)\,
\exp\!\left[2\pi i(hx+ky+lz)\right]dV ,
\label{eq:ch04b-FT}
\end{equation}
反向变换给出 $\rho(x,y,z)=V_{c}^{-1}\sum_{h}\sum_{k}\sum_{l}F_{hkl}
\exp\!\left[-2\pi i(hx+ky+lz)\right]$。这一对变换关系是结构分析的基础，也解释了一个根本困难：
实测强度只给出 $\left|F_{hkl}\right|^{2}$，相角 $\alpha$ 在测量中丢失，因此“从衍射数据解结构”
的实质是先用 Patterson 法、直接法或全谱拟合（Rietveld）解决相角问题。

此处已对照原页面图像（书内 p.250 / PDF p.269）：原书正变换印为
$F_{hkl}=\int_{V}\rho(x,y,z)\exp[2\pi i(hx+ky+lz)]\,dV$，反向变换印为
$\rho(x,y,z)=\frac{1}{V}\sum_{h}\sum_{k}\sum_{l}F_{hkl}\exp[-2\pi i(hx+ky+lz)]$，
即积分在晶胞体积 $V$ 上取、微分元为 $dV$，归一化常数 $1/V$ 只出现在反向变换一侧。本式按此
改写——原稿正变换写成单位立方上的 $\int_{0}^{1}\!\!\int_{0}^{1}\!\!\int_{0}^{1}\cdots dx\,dy\,dz$，
与反向变换的 $V_{c}^{-1}$ 不属同一套归一化约定，已改成一致写法。原书同页另印出 Friedel 定律的
推导（由 $I_{hkl}=\left|F_{hkl}\right|^{2}=F_{hkl}F_{hkl}^{*}$ 与 $I_{\bar{h}\bar{k}\bar{l}}$
逐项展开相等得 $\left|F_{hkl}\right|=\left|F_{\bar{h}\bar{k}\bar{l}}\right|$），并指出衍射空间
强度分布的中心对称对应 $11$ 种 Laue 群。

\subsubsection{衍射的系统消光}

在 §4.3 的衍射几何讨论中，衍射条件由倒易格子点与 Ewald 反射球相交给出。若正空间格子为复格子，
倒易格子点会系统消失；晶体中存在螺旋轴与滑移面等微观对称元素时，也会引起衍射的系统缺失。
用式~\equref{eq:ch04b-F} 直接计算即可得到全部消光规律。

\paragraph{体心格子}晶胞中任一原子位于 $(x,y,z)$ 时，必有同种原子位于
$(x+\tfrac12,y+\tfrac12,z+\tfrac12)$。代入结构因子并提取公因子：
\begin{equation}
F_{hkl}=\sum_{j}f_{j}\exp\!\left[2\pi i(hx_{j}+ky_{j}+lz_{j})\right]
\left\{1+\exp\!\left[\pi i(h+k+l)\right]\right\}.
\label{eq:ch04b-I-extinct}
\end{equation}
当 $h+k+l=2n$（偶数）时括号内为 $2$，$F_{hkl}=2\sum_{j}f_{j}$，衍射出现；当 $h+k+l=2n+1$
（奇数）时括号内为零，$F_{hkl}=0$，衍射不出现。即体心格子的倒易格子中 $h+k+l$ 为奇数的点
全部消失——$100$、$010$、$001$ 不出现，出现的是 $200$、$020$、$002$、$110$、$101$、$011$。
该倒易格子是边长为 $2a^{*},2b^{*},2c^{*}$ 的\textbf{面心}格子。

\paragraph{面心格子}晶胞中任一原子位于 $(x,y,z)$ 时，必有另外三个原子分别位于
$(x+\tfrac12,y+\tfrac12,z)$、$(x+\tfrac12,y,z+\tfrac12)$、$(x,y+\tfrac12,z+\tfrac12)$。代入后
\begin{equation}
F_{hkl}=\sum_{j}f_{j}\exp\!\left[2\pi i(hx_{j}+ky_{j}+lz_{j})\right]
\left\{1+e^{\pi i(h+k)}+e^{\pi i(h+l)}+e^{\pi i(k+l)}\right\}.
\label{eq:ch04b-F-extinct}
\end{equation}
$h,k,l$ 全为奇数或全为偶数时 $h+k$、$h+l$、$k+l$ 均为偶数，括号内为 $4$，
$F_{hkl}=4\sum_{j}f_{j}\exp\!\left[2\pi i(hx_{j}+ky_{j}+lz_{j})\right]$，衍射出现；$h,k,l$
奇偶混杂时三个指数中一个为偶数、两个为奇数，括号内为零，衍射不出现。即面心格子的倒易格子
中 $100$、$010$、$001$、$110$、$101$ 不出现，而 $200$、$020$、$002$、$111$ 出现；其倒易格子
的基元是体心格子。类似地可证底心格子的倒易格子仍是底心格子。

\paragraph{滑移面与螺旋轴}设垂直于 $c$ 轴存在 $b$ 滑移面（设滑移面过 $z=0$），晶胞中任一原子
$(x,y,z)$ 必对应 $(x,y+\tfrac12,-z)$，代入后对 $hk0$ 指标可消去 $z$ 的影响，得到
$F_{hk0}=\sum_{j}f_{j}\exp\!\left[2\pi i(hx_{j}+ky_{j})\right]\left[1+\cos(\pi k)\right]$
（原书此处印 $1+\exp(k\pi i)$）。因此 $k=2n$ 时括号内为 $2$，衍射出现；$k=2n+1$ 时括号内为
零，$F_{hk0}=0$：倒易格子中 $hk0$ 面上 $k$ 为奇数的点全部消失。
同理，平行于 $b$ 方向存在 $2_{1}$ 螺旋轴时，$0k0$ 衍射中 $k$ 为奇数的点消失；垂直于 $b$
方向存在 $a$ 滑移面时，$(h0l)$ 中 $h=2n+1$ 的点消失；平行于 $c$ 存在 $4_{1}$ 螺旋轴时，
$00l$ 中只有 $l=4n$ 的点存在。此处已对照原页面图像（书内 p.252—253 / PDF p.271—272）：
滑移面一段原书取原子对 $(x,y+\tfrac12,-z)$（滑移面过 $z=0$），$2_{1}$ 螺旋轴一段原书取
$(-x,y+\tfrac12,-z)$（螺旋轴过 $x=0,z=0$），结论均为对应指标中奇数的点消失。原书在同节
末尾以三条结论归纳：P 格子且无螺旋轴、滑移面等微观对称元素时衍射谱无任何系统消光；复格子
使某些 $hkl$ 系统消光但不破坏倒易格子特性；螺旋轴与滑移面使倒易格子出现「空位缺陷」，
并明确举出垂直于 $b$ 的 $a$ 滑移面使 $(h0l)$ 中 $h=2n+1$ 消失、平行于 $c$ 的 $4_{1}$
螺旋轴使 $00l$ 中仅 $l=4n$ 存在。

\begin{table}[tbp]
  \centering
  \caption{点阵类型与微观对称元素引起的系统消光规律。表中“出现条件”是对该组衍射指标而言，
    其余指标不受此条限制。\src{据 S1（上册）§4.4.4，书内 p.250—253 / PDF p.269—272 归纳。}}
  \label{tab:ch04b-extinction}
  \begin{tabular}{@{}lll@{}}
    \toprule
    类型 & 出现（衍射）条件 & 系统消光 \\
    \midrule
    简单格子 $P$ & 无附加条件 & 无消光，全部 $hkl$ 出现 \\
    体心 $I$ & $h+k+l=2n$ & $h+k+l$ 为奇数 \\
    面心 $F$ & $h,k,l$ 全奇或全偶 & $h,k,l$ 奇偶混杂 \\
    底心 $C$ & $h+k=2n$ & $h+k$ 为奇数 \\
    $b$ 滑移面 $\perp c$ & $hk0$ 中 $k=2n$ & $hk0$ 中 $k$ 为奇数 \\
    $a$ 滑移面 $\perp b$ & $h0l$ 中 $h=2n$ & $h0l$ 中 $h$ 为奇数 \\
    $2_{1}$ 螺旋轴 $\parallel b$ & $0k0$ 中 $k=2n$ & $0k0$ 中 $k$ 为奇数 \\
    $4_{1}$ 螺旋轴 $\parallel c$ & $00l$ 中 $l=4n$ & $00l$ 中 $l\neq4n$ \\
    \bottomrule
  \end{tabular}
\end{table}

\begin{keycond}
三类信息可以直接由“哪些峰存在”读出：点阵类型（复格子）使\textbf{整类}衍射指标消失，倒易
格子的格子型式甚至基本矢量发生改变而不破坏格子特性；螺旋轴与滑移面使\textbf{某条直线或某个面}
上的格子点部分消失，相当于在倒易格子里出现“空位缺陷”。因此，衍射谱上峰的“有无规律”是点阵
类型与可能空间群的直接证据，与峰位、峰强构成相互独立的三重校验。反过来，若一套指标化结果
不能给出任何系统消光规律，其可靠性就要打折扣（见 §4.6.4）。
\end{keycond}

\begin{methodsum}
本章主线可以压缩成三句话：\textbf{峰位 $\rightarrow$ 晶胞}（由 $2\theta$ 得 $d$、$Q$，解出晶胞
参数与点阵对称性）；\textbf{峰强 $\rightarrow$ 晶胞内原子分布}（由修正后的 $\left|F_{hkl}\right|^{2}$
得到原子种类、坐标、占位度与热振动）；\textbf{峰的有无 $\rightarrow$ 点阵类型与微观对称元素}
（系统消光定格子类型、滑移面与螺旋轴）。三者的关系是：布拉格方程给出衍射出现的必要条件，
结构因子不为零才是充分条件。\src{书内 p.253—254 / PDF p.272—273}
\end{methodsum}

\subsection{物相鉴定与相定量分析}
\src{上册 §4.5，书内 p.254—264 / PDF p.273—283}

\subsubsection{物相定性分析}

\begin{probout}
\textbf{测什么}：把实测 $d$—$I/I_{1}$ 数据组与标准卡片比对，确定样品由哪些晶相组成。\\
\textbf{能推断什么}：某物相\textbf{存在}（当其特征线全部出现且相对强度合理时），以及通过
与已知结构对比给出晶型是否相同的判断。\\
\textbf{不能推断什么}：某物相\textbf{不存在}。含量低或衍射能力弱时其特征线显示不出来，
因此该方法只能确定某相存在，不能确定某相绝对不存在；含量低于约 $3\%$ 的物相通常很难鉴定出来，
若某物相全部由轻原子（C、H、O、N）组成而其他相含重原子，即使含量到 $40\%$ 也可能鉴别不出来。
\src{书内 p.258 / PDF p.277}
\end{probout}

\begin{principle}
物相定性的依据有两条。第一，任何一种结晶固体都有自己独立的一套衍射图谱，峰位与强度完全
取决于其内部结构：由式~\equref{eq:ch04b-Q}，晶胞参数不同则 $d$ 不同、衍射角不同；由
式~\equref{eq:ch04b-F}，各峰强度由晶胞中原子分布方式决定。因此每种晶态物质的衍射图谱与其
结构之间存在一一对应关系。第二，物质不会因为与其他物质混合而改变自己的衍射行为——混合物中
各物相的衍射互不干扰，混合物的图谱是各组成物相图谱的简单叠加。这两条合起来就是“用图谱
鉴别晶态物质”的全部依据。
\end{principle}

\begin{measure}
标准数据的形式与检索方式决定了物相分析的工作流程。1938 年 Hanawalt 等人用 $d$—$I$ 数据组
代替衍射图制作卡片；1942 年美国测试与材料协会将数据汇编成 1300 张卡片出版，简称 ASTM 卡片；
1972 年后由粉末衍射标准联合委员会（JCPDS）编辑出版，称 JCPDS 卡（PDF 卡）；现在由国际衍射
数据中心（ICDD）发行，简称 ICDD 卡。卡片按区域分载：卡片号（组号 + 卡号）、该相三强线及第一个
峰的 $d$ 值与相对强度、化学式与英文名称、数据可靠度符号（用星号标记者表示数据非常可靠，
$i$ 表示比较可靠，空白表示可靠性一般，$O$ 表示可靠性差，$C$ 表示由晶体结构数据计算得到，个别卡片为 $R$ 表示来自 Rietveld 精修结果）、
实验条件（辐射种类 $\mathrm{Rad}$、波长 $\lambda$、滤波片 $\mathrm{Filter}$、相机直径
$\mathrm{Dia}$、测量方法、参考文献）、晶体学数据（晶系 $\mathrm{sys}$、空间群 $\mathrm{S.G}$、
晶胞参数、计算密度 $D_{x}$、晶胞中化学式单位数 $Z$）、物性数据、以及全部 $d$—$I$—$hkl$ 数据
（强度以最强峰 $I_{1}=100$ 归一）。\\
检索工具有两类：\textbf{字母索引}（Davey 索引，按物相英文名称字母顺序，条目含名称、化学式、
三强线 $d$ 值与相对强度、卡片号、参比强度）最常用，适用于已知化学组成、能估计可能物相的样品；
\textbf{数字索引}（Hanawalt 索引，按最强线的 $d$ 值由大到小分 87 组，组内按第二、第三个 $d$ 值
依次排列，且前三条线的 $d$ 值轮流占据首位，因此同一张卡片在索引中出现三次）适用于物相或元素
完全未知的样品。
\end{measure}

\begin{dataflow}
（1）按要求制样并获得衍射数据；（2）由峰位按布拉格方程计算 $d$ 值，读取强度并计算相对强度
$I/I_{1}$；（3）用字母或数字索引检索 PDF 卡片；（4）与候选卡片逐线比对，判定唯一合理的 PDF
卡片。比对时允许实测 $d$ 与卡片 $d$ 存在差异，原书举例如下：$d\sim1$ Å 时允许 $\pm0.003$ Å、
$d\sim3.5$ Å 时允许 $\pm0.03$ Å、$d\sim8$ Å 时允许 $\pm0.25$ Å；误差在此范围内即视为相符。
（此处已对照原页面图像（书内 p.257 / PDF p.276）：原书三档误差都给出了对应的 $d$ 值，即
$1$ Å、$3.5$ Å、$8$ Å；本节此前只写“更高角度允许 $\pm0.25$ Å”，未指认其适用的 $d$ 范围。）
实验时要注意温度与仪器零点的影响，消除仪器误差可加内标物（如 \ce{SiO2}、\ce{CaF2}）。
\end{dataflow}

\begin{pitfall}
\textbf{（1）把 $d$ 值符合与强度符合的权重搞反。}实测数据与 PDF 卡片通常不完全一致；
比对时 $d$ 值的符合比相对强度的符合更重要，相对强度只作参考。\\
\textbf{（2）忽略低角区与高角区的信息量差异。}不同晶体在低角度 $d$ 值一致的机会很少，而在高角度
不同晶体间衍射峰相似的机会较大。因此低角区衍射与卡片数据的符合比高角区的符合更重要。\\
\textbf{（3）把重叠峰当成某相的强线。}多相混合样品中不同相的某些衍射峰可能互相重叠，某些“强线”
实际上并不是某一物质的强衍射；若以它作为最强线去检索，很难得到符合的结果。\\
\textbf{（4）忽略仪器分辨率与样品位移对“匹配”的影响。}不同仪器（狭缝、测角仪半径、光源）给出的
峰宽与峰形不同，样品位移还会使全部峰位同向偏移。若把 PDF 卡片当成“精确到小数点后三位就必须
逐位对上”的标尺，会做出错误排除或错误匹配。正确做法：先用内标校正峰位，再按 $d$ 值的角度分段
容差比对，并优先信任低角线。\\
\textbf{（5）结构相似相之间的误判。}有些物质结构相似，仅点阵常数差别不大、原子散射能力也相近，
衍射峰差别很小，必须与化学分析、电子探针、能谱等方法结合才能得出正确结果。\\
\textbf{（6）忽略卡片版本。}不同编号的同一物质卡片，应以发表较晚卡片上的数据为准。
\end{pitfall}

\begin{crosscheck}
物相鉴定不应单独成立：应与化学分析、X 射线荧光（元素组成）、光谱、电子探针与能谱（微区成分）
互相配合。若物相太多、衍射线重叠严重，可用重力、磁力等方法先把某些物相富集或分离，再分别鉴定。
鉴定结论还应与样品来源或制备方法一致——例如由硝酸盐分解制备的样品中出现对应的氧化物相是合理的，
出现无关的硅酸盐相则需要解释来源。
\end{crosscheck}

\subsubsection{物相定量分析}

\begin{probout}
\textbf{测什么}：各晶相在样品中的相对含量（质量分数）。\\
\textbf{前提}：物相定性分析是相定量的出发点——只有先确认样品中含有指定物相，相定量才有意义。\\
\textbf{能推断什么}：晶相之间的相对比例，精度通常到百分之一到百分之几。\\
\textbf{不能推断什么}：不能直接给出非晶相的含量，也不能在未做标定的情况下给出绝对含量。
\end{probout}

\begin{principle}
由式~\equref{eq:ch04b-intensity}，对选定 $hkl$ 的衍射峰，$\left|F_{hkl}\right|^{2}$、$j$、$D$
以及与角度相关的 $P$、$L$ 都是确定值，强度仅是样品受照射有效体积 $\Delta V$ 的函数，可把确定
的项合并为常数：$I=K\Delta V$。采用 Bragg--Brentano 衍射仪、平板样品时
$\Delta V=A_{0}/(2\mu_{l})$，故 $I=KA_{0}/(2\mu_{l})$，其中 $\mu_{l}$ 为线吸收系数，
$A_{0}$ 为入射线垂直截面积。\textbf{衍射仪做相定量的最大便利正在于此：吸收效应不随 $\theta$
改变，易于校正。}\\
多相体系中各相独立产生衍射，混合体系可看作各单相体系的体积权重加和。设 $i$ 相的体积分数为
$v_{i}$、质量分数为 $w_{i}$，总线吸收系数 $\mu_{l}=\sum_{i}v_{i}\mu_{l,i}$，则
\begin{equation}
I_{i}=\frac{K_{i}v_{i}}{\mu_{l}} .
\label{eq:ch04b-Ii-V}
\end{equation}
改用质量吸收系数 $\mu_{m}$（$\mu_{l}=\rho\mu_{m}$，$\rho$ 为样品密度；$\mu_{m}$ 只与组成元素及
其质量分数有关，与聚集状态无关，故同分异构体的多相体系 $\mu_{m}$ 是常数）：
\begin{equation}
\mu_{m}=\sum_{i}w_{i}\mu_{m,i},\qquad
I_{i}=\frac{k_{i}w_{i}}{\mu_{m}},\qquad k_{i}=\frac{K_{i}}{\rho_{i}} .
\label{eq:ch04b-Ii-w}
\end{equation}
式~\equref{eq:ch04b-Ii-w} 是 X 射线相定量分析最基本的公式，以下各种方法都从它出发；各方法的
差别只在于如何计算或消去常数 $k_{i}$ 与总质量吸收系数 $\mu_{m}$。
（此处已对照原页面图像（书内 p.259—260 / PDF p.278—279）：原书由单相式 $I=K\Delta V$ 与
Bragg-Brentano 有效体积 $\Delta V=A_{0}/(2\mu)$ 得 $I_{i}=K_{i}v_{i}/\mu_{t}$（$\mu_{t}$ 为总线
吸收系数），再由 $v_{i}=V_{i}/V_{t}=w_{i}\mu_{t}/(\rho_{i}\mu_{mt})$ 过渡到质量分数，得
\[ I_{i}=\frac{K_{i}}{\rho_{i}}\cdot\frac{w_{i}}{\mu_{mt}}, \]
即本式——$\rho_{i}$ 与 $K_{i}$ 同处分子，故 $k_{i}=K_{i}/\rho_{i}$ 确为合并写法。符号差异：原书总线
吸收系数与总质量吸收系数记作 $\mu_{t}$、$\mu_{mt}$，本节沿用 $\mu_{l}$、$\mu_{m}$；原书未印
$\mu_{m}=\sum_{i}w_{i}\mu_{m,i}$ 的质量加权求和式，只印 $\mu_{mt}\rho_{t}=\mu_{t}=\sum_{i}v_{i}\mu_{i}$，
两者等价。）
\end{principle}

\paragraph{外标法}若多相体系中各相化学组成相同而结构不同（同分异构体混合物），则 $\mu_{m}$ 是
常数、不随各组分质量分数改变，于是相同实验条件下纯 $i$ 相该峰的强度 $I_{i}^{0}$ 与混合样品中
的强度之比直接给出含量：
\begin{equation}
w_{i}=\frac{I_{i}}{I_{i}^{0}} .
\label{eq:ch04b-external}
\end{equation}
外标法原理简单，但局限性大：只适用于上述有限体系（$\mu_{m}$ 恒定），需要纯物质作外标物，
且必须严格控制实验条件——制样条件（样品处理、压样用力）平行、仪器稳定，通常需要多次制样与
测试以检验强度的重复性。

\paragraph{内标法}对 $n>2$ 相体系，各组分 $\mu_{m,i}$ 不同，$\mu_{m}$ 随组成变化，
式~\eqref{eq:ch04b-external} 不再成立。此时在样品中定量加入一种原体系不含的标准物质 $s$ 相。
设原体系在新混合体系中的质量分数为 $w_{0}$，$i$ 相与原体系的质量分数分别为 $w_{i}$、$w_{s}$，
则 $i$ 相在新体系中的质量分数 $w_{i}'=w_{i}w_{0}$。将 $i$ 相与 $s$ 相的强度公式相除，总质量
吸收系数被消去：
\begin{equation}
\frac{I_{i}}{I_{s}}=K\frac{w_{i}'}{w_{s}},\qquad w_{i}=\frac{w_{i}'}{w_{0}}=\frac{w_{i}'}{1-w_{s}} .
\label{eq:ch04b-internal}
\end{equation}
即 $i$ 相与内标相的强度比与其在新混合体系中的含量成正比。实际工作中配制一系列 $i$ 与 $s$ 含量
已知的样品，测强度比作工作曲线，由斜率求 $K$，再算出 $w_{i}'$ 并推回原体系的 $w_{i}$。
内标物的选择条件：测试峰与待测组分峰尽量靠近但不重叠、也不与其他峰重叠；峰强度足够高；
结晶完整性高、无择优取向（一般优先选立方晶系物质）；性质稳定，X 射线照射下不变、不与样品
起化学反应。内标法所得校正曲线适用于任何含待测物相的多相体系，只要求用相同内标物与相同衍射峰；
若校正曲线所用峰在样品中与其他组分重叠，可另选新峰，通过新峰与原峰积分强度之比对各点乘以该
比值即可换算曲线，不必重做实验。

\paragraph{参比强度法（RIR／$K$ 值法）}参比强度法是内标法的发展，本质仍是内标法，区别在于
\textbf{统一了参比物}：以刚玉（corundum，$\alpha$-\ce{Al2O3}，$R\bar{3}c$）为通用参考标准，
取刚玉 $(113)$ 衍射峰为参比峰，定义待测相 $i$ 与刚玉按质量比 $1:1$ 混合时二相最强线的强度比
为该相的参比强度 $K_{i}$。当体系中原有的 $n$ 相之外加入定量内标物 $f$ 时，
\begin{equation}
\frac{I_{i}}{I_{f}}=\frac{K_{i}}{K_{f}}\cdot\frac{w_{i}}{w_{f}},
\label{eq:ch04b-RIR}
\end{equation}
即体系中任意二组分的强度比正比于二组分的浓度比，比例系数等于二组分参比强度之比，\textbf{与样品中
其他组分的存在无关}。因此只要在同一扫描中测出待测相峰与内标物峰的强度比，配合参比强度与内标物
加入量即可算出含量，\textbf{不必再做校正曲线}——内标物在这里起的是“冲洗剂”的作用，故也称基体
冲洗法。1972 年之后 JCPDS 卡片直接给出部分物质与刚玉相比的参比强度值；使用卡片值时必须注意
仪器条件与样品状况，必要时自行测定。若以刚玉为内标物加入 $n$ 组分体系，且令 $K_{f}=K_{c}=1$，
公式进一步简化为 $I_{i}/I_{c}=K_{i}w_{i}/w_{c}$，即 $w_{i}=w_{c}I_{i}/(K_{i}I_{c})$。

\paragraph{自动冲洗绝热法（无标法）}在 $n$ 相体系中不加内标物，仅从一张衍射图求各相含量。
其原理与参比强度法相同：任意两相之间 $I_{i}/I_{j}=(K_{i}w_{i})/(K_{j}w_{j})$，对 $n$ 相体系可写
出 $n(n-1)/2$ 个方程，其中独立方程数为 $n-1$，加上归一化条件 $\sum_{i}w_{i}=1$ 正好有 $n$ 个
独立方程。由此得
\begin{equation}
w_{i}=\frac{I_{i}/K_{i}}{\sum_{j=1}^{n}I_{j}/K_{j}} .
\label{eq:ch04b-flushing}
\end{equation}
若混合样中各组分都是可结晶物质且都有已知的 $K_{i}$，则无需加入冲洗剂即可直接计算各组分含量。
\textbf{但当样品中存在无定形（非晶）物质时，该公式不能应用}，只能向样品中加入内标物处理——
原因很简单：非晶相不产生布拉格峰，$\sum w_{i}=1$ 的归一化假设被破坏，算出的各晶相含量会被系统性
高估。
（此处已对照原页面图像（书内 p.263 / PDF p.282）：原书由参比强度法关系
$I_{1}/I_{i}=(K_{1}/K_{i})(w_{1}/w_{i})$（$i=1,2,\cdots,n$）出发，指出 $i=1$ 的方程是恒等式、非独立
方程；改写为 $w_{i}=(I_{i}/K_{i})(K_{1}/I_{1})w_{1}$ 后代入归一化方程 $\sum_{i}w_{i}=1$，得
$w_{1}=\left(\frac{K_{1}}{I_{1}}\sum_{i=1}^{n}\frac{I_{i}}{K_{i}}\right)^{-1}$，对体系中任意一相 $i$ 即
$w_{i}=\left(\frac{K_{i}}{I_{i}}\sum_{j=1}^{n}\frac{I_{j}}{K_{j}}\right)^{-1}$，与式~\equref{eq:ch04b-flushing}
互为倒数写法、代数等价。原书同页明确「原理上，自动冲洗绝热法与参比强度法相同」，各组分「均有已知的
参比强度 $K_{i}$ 值」，故 $K_{i}$ 即参比强度；「当样品中有无定形物质存在时，就不能应用这个公式来计算」
与本段末句一致。）

\begin{table}[tbp]
  \centering
  \caption{四类相定量方法的适用差别。选择方法的第一步是判断“总质量吸收系数是否随组成变化”，
    第二步是判断“能否接受加入内标物”。\src{据 S1（上册）§4.5.2，书内 p.258—264 / PDF p.277—283 归纳。}}
  \label{tab:ch04b-quantmethods}
  \begin{tabular}{@{}p{21mm}p{42mm}p{33mm}p{38mm}@{}}
    \toprule
    方法 & 适用体系与前提 & 标样与工作曲线 & 主要误差来源 \\
    \midrule
    外标法 & 各相化学组成相同、结构不同（同分异构体）；$\mu_{m}$ 恒定 &
    需纯 $i$ 相作外标；需严格控制制样与仪器条件，常需多次重复 &
    制样平行性、压样用力、仪器漂移；适用范围窄 \\
    内标法 & 任意 $n$ 相体系（$n>2$）；可加入原体系不含的 $s$ 相 &
    需选合格内标物；需配制系列样品作工作曲线 &
    内标物纯度与加入量准确性、混合均匀性、峰重叠 \\
    参比强度法（RIR／$K$ 值） & 任意多相体系；本质是内标法 &
    以刚玉 $(113)$ 峰为统一参比；可用卡片 $K_{i}$ 值，\textbf{不必做工作曲线} &
    卡片 $K_{i}$ 值与本机条件不匹配、择优取向、峰重叠 \\
    自动冲洗绝热法 & 各相均可结晶、均有已知 $K_{i}$ 的多相体系；\textbf{不能含非晶相} &
    无需内标物、无需工作曲线，仅用一张衍射图 &
    存在非晶相时归一化失效、$K_{i}$ 值误差直接传递、重叠峰 \\
    \bottomrule
  \end{tabular}
\end{table}

\begin{measure}
相定量的实验要点：（1）样品尽可能细，粒度通常要小于 $15$\,µm，并尽可能消除择优取向；
（2）合理选择衍射条件——管压、管流、狭缝、扫描速度、扫描区间，数据处理要审慎，必要时重复测试；
（3）\textbf{峰强用积分强度且计数要足够}；若用峰高代替积分强度，要注意仪器型号与实验条件——
同型号仪器的峰高／峰宽比随 $2\theta$ 的变化相同，不同型号仪器则不同，因此用峰高得到的校正曲线
只能在同一型号仪器上通用，用积分强度得到的校正曲线则可在不同型号仪器上通用；若待测组分晶粒
尺寸在 $100$ nm 以下且各样品的晶粒大小不一致，则不能用峰高代替积分强度；（4）相定量的灵敏度
约为 $5\%$：含量较高时一般可确认该物相存在，含量很低或无法给出时，需要结合样品来源或制备
方法、化学分析及其他表征数据综合判断。\src{书内 p.264 / PDF p.283}
\end{measure}

\begin{pitfall}
\textbf{定量误差主要来自三处。}\\
\textbf{（1）择优取向（最常被低估的一项）。}层状、针状、片状晶粒在压片时趋于定向排列，使
反射面平行于样品面的衍射峰被增强、垂直的被减弱（图~\figref{fig:ch04b-texture}）。由于各相的
取向程度不同，相与相之间的强度比不再正比于含量比，误差可达数十个百分点。抑制办法：减小粒度、
侧装或轻压、加冲淡剂（如与样品吸收系数接近的稀释剂）、必要时用取向函数在 Rietveld 精修中校正。\\
\textbf{（2）微吸收。}颗粒尺寸与吸收系数不匹配时，X 射线在单个颗粒内被吸收，颗粒内部的散射
被屏蔽，导致“强度—含量”偏离线性。颗粒越粗、各相吸收系数差别越大，偏差越严重；吸收系数大
的相（含重原子）在粗颗粒状态下其强度会被系统性压低。判据是颗粒线度应远小于 $1/\mu_{l}$；
否则需减小粒度或按 Brindley 方法做微吸收校正。\\
\textbf{（3）非晶相未计入。}自动冲洗绝热法（以及一切用 $\sum w_{i}=1$ 归一的方法）以“全部
物质都产生布拉格峰”为前提。催化剂、玻璃、无定形碳、无定形载体上负载的活性相都含有非晶成分，
此时归一化把非晶部分的质量摊给了各晶相，\textbf{各晶相含量被系统性高估}，而且高估比例随非晶
含量增大而增大。正确做法是加入内标物，由内标相的含量反推非晶相的总量，再换算各晶相在原样品
中的绝对含量。\\
此外还有两项次要但常被忽略的来源：峰重叠未分离（把重叠峰强度全部计入某一相）、以及用峰高
代替积分强度（图~\ref{fig:ch04b-pld} 所示的角度依赖使峰宽随 $2\theta$ 变化，峰高比不再正比于
积分强度）。
\end{pitfall}

\begin{crosscheck}
相定量结果应与独立方法互校：热分析（灼烧失重、结晶水）、元素分析（化学法或 XRF）给出元素
总量约束；Rietveld 全谱拟合给出包含结构信息的相含量（见后续分片）；把内标法与参比强度法
并行做一次，若两者相差超过预期误差，问题通常出在择优取向或非晶相未计入，而不是在计算方法上。
\end{crosscheck}

\subsection{衍射图的指标化}
\src{上册 §4.6，书内 p.264—269 / PDF p.283—288}

\begin{probout}
\textbf{测什么}：为多晶衍射图中的每一条衍射线指定衍射指标 $hkl$。\\
\textbf{能推断什么}：由
\begin{equation}
Q_{hkl}=f\left(h,k,l;a,b,c,\alpha,\beta,\gamma\right)
\label{eq:ch04b-indexbase}
\end{equation}
反解出晶胞参数 $(a,b,c,\alpha,\beta,\gamma)$，实现倒易点阵的重建；同时得到晶体所属晶系、点阵
类型与可能的空间群等信息。用一套晶胞参数指标化全部衍射线，是确认样品为单相的重要依据；
把指标化所得的晶胞参数、样品实际组成与已知结构的原子参数结合起来计算衍射图并与实验图对照，
可以判明样品是否与已知化合物具有相同晶型。\\
\textbf{困难在哪}：$n$ 条线只有 $n$ 个方程，一般晶系含 $3n+6$ 个未知数（即使对称性最高的立方
晶系也有 $3n+1$ 个），从数学上看指标化是一个“无解”的难题；出路在于衍射指标只能取整数，
且低角度区域首先出现指标低的衍射峰——指标化因此是尝试与推理结合的过程，本质上是\textbf{多解}
问题。\src{书内 p.264—265 / PDF p.283—284}
\end{probout}

\subsubsection{立方晶系指标化方法：解析法}

立方晶系只有一个晶胞参数需要确定，不同衍射峰的 $Q$ 值之间存在完全由整数指标决定的比例关系：
\begin{equation}
\frac{Q_{i}}{Q_{1}}=\frac{h_{i}^{2}+k_{i}^{2}+l_{i}^{2}}{h_{1}^{2}+k_{1}^{2}+l_{1}^{2}} .
\label{eq:ch04b-cubicratio}
\end{equation}
由于指标只能取整数且以平方和的形式出现，只需按 $0,1,2,3,\dots$ 从小到大组合 $h,k,l$，依次得到
$100$、$110$、$111$、$200$、$210$、$211$、$220$、$300$、$310$、$311$、$222$、$320$、
$321$、$400$、$410$、$322$、$411$、$330$、$331$、$420$ 等，再按系统消光规律剔除不出现的指标，
即可得到 $Q$ 比值序列：

（1）\textbf{简单立方}：$Q$ 比值 $1:2:3:4:5:6:8:9:10:11:12:13:14:16:17:18\cdots$，
其中缺失 $7$ 和 $15$（$h^{2}+k^{2}+l^{2}$ 不能等于 $7$ 或 $15$）。\\
（2）\textbf{体心立方}：只有 $h+k+l=2n$ 的衍射出现，$Q$ 比值
$2:4:6:8:10:12:14:16:18:20\cdots$，化为最简后即 $1:2:3:4:5:6:7:8:9:10\cdots$，\textbf{所有整数
都出现}（比值 $7$ 出现，这是与简单立方的关键差别）。\\
（3）\textbf{面心立方}：$h,k,l$ 奇偶混杂的衍射消光，$Q$ 比值
$3:4:8:11:12:16:19:20:24:27\cdots$，呈现疏密交替的衍射花样——这也是面心立方粉末图最直观的
识别特征。（此处已对照原页面图像（书内 p.265 / PDF p.284）：原书三式分别只印到
$1:2:3:4:5:6:8:9:10:11:12:13:14:16:17:18\cdots$（并注明缺失 $7$ 和 $15$）、
$2:4:6:8:10:12:14:16:18:20\cdots$（化简后为 $1:2:3:4:5:6:7:8:9:10\cdots$，并注明其中出现 $7$）
与 $3:4:8:11:12:16:19:20\cdots$，三式均以省略号收尾；本序列续列的 $24$、$27$ 等项是按
$h^{2}+k^{2}+l^{2}$ 的可能取值与系统消光条件补出的，不是原书印出的项。）

\begin{keycond}
立方晶系多晶衍射数据指标化的一般步骤（原书五步，此处给出每步的判据与常见失误）：
（1）读出各衍射峰的 $2\theta$，由布拉格方程 $\sin\theta=\lambda/(2d)$ 计算 $\sin^{2}\theta$
（同一辐射源下 $\sin^{2}\theta$ 的比值与 $Q$ 的比值相同），按由小到大排列；
（2）计算各 $\sin^{2}\theta$ 与最小 $\sin^{2}\theta$ 的比值；
（3）若比值都接近整数则化为整数；若有的比值不接近整数，观察数值关系乘以适当的整数（$2$、$3$、
$\dots$）使全部比值接近整数并进一步化简——这一步是立方解析法的关键，也是产生“晶胞体积
亚倍／倍数假解”的地方；
（4）根据比值关系判断格子类型，确定各衍射峰的指标 $hkl$；
（5）由布拉格方程计算各峰的 $d$ 值，进而计算晶胞参数 $a$。
\textbf{判据}：第（4）步得到的格子类型必须与该相已知的消光规律（表~\tabref{tab:ch04b-extinction}）
一致；第（5）步中各条线单独算出的 $a$ 应当在误差范围内一致，若 $a$ 随角度系统漂移，说明峰位
未做位移校正或指标化有误。
\end{keycond}

\begin{example}[面心立方相的指标化（模型数据）]
以下为一组自洽的模型数据（$a=4.078$ Å，Cu $K\alpha$，$\lambda=1.5418$ Å，由
$\sin^{2}\theta=(\lambda^{2}/4a^{2})N$ 正算得到），用来演示上节五步：
\begin{center}
\begin{tabular}{@{}ccccc@{}}
  \toprule
  $2\theta$ / (°) & $\sin^{2}\theta$ & 比值（$\times3$） & $N=h^{2}+k^{2}+l^{2}$ & $hkl$ \\
  \midrule
  38.22 & 0.1072 & 3 & 3 & 111 \\
  44.42 & 0.1429 & 4 & 4 & 200 \\
  64.63 & 0.2859 & 8 & 8 & 220 \\
  77.66 & 0.3931 & 11 & 11 & 311 \\
  81.80 & 0.4288 & 12 & 12 & 222 \\
  98.24 & 0.5718 & 16 & 16 & 400 \\
  111.02 & 0.6790 & 19 & 19 & 331 \\
  115.40 & 0.7147 & 20 & 20 & 420 \\
  \bottomrule
\end{tabular}
\end{center}
\vspace{0.3em}
（1）由 $2\theta$ 算出 $\sin^{2}\theta$ 并升序排列；（2）与最小值 $0.1072$ 相比得
$1:1.333:2.667:3.667:4:5.333:6.333:6.667$，都不是整数；（3）观察数值关系，乘以 $3$ 后得
$3:4:8:11:12:16:19:20$——全部为整数，且不含 $3:1$ 的简单比例（说明最小线不是 $100$），
序列表中缺失 $5,6,7,9,10$ 等值；（4）该序列正是面心立方的特征序列（最小 $N=3$ 对应 $111$），
据此判定格子类型为面心立方并指定上表指标；（5）由
$a=\lambda\sqrt{N}/(2\sin\theta)$ 逐线计算，八条线一致给出 $a=4.078$ Å（高角线权重更大，
因为 $\Delta d/d$ 随 $\theta$ 增大而减小）。
\textbf{本例的用途}：它演示了“比值 $\rightarrow$ 整数化 $\rightarrow$ 消光判据 $\rightarrow$ 指标”
这一完整链条；若在第（3）步只乘到 $1$ 而不乘 $3$，就会得到 $1:1.33:2.67\cdots$ 的“非整数序列”
而误判为不合格数据。
\end{example}

\subsubsection{Hesse--Lipson 解析法}

\begin{principle}
Hesse--Lipson 法适用于四方、六方与正交晶系，思路仍是利用晶体周期性结构决定的衍射峰之间的
变化关系。以正交晶系（三个待定参数）为例，把布拉格方程代入 $\sin^{2}\theta$ 得
\begin{equation}
q_{hkl}\equiv\sin^{2}\theta_{hkl}
=\frac{\lambda^{2}}{4}\left(\frac{h^{2}}{a^{2}}+\frac{k^{2}}{b^{2}}+\frac{l^{2}}{c^{2}}\right)
=A h^{2}+B k^{2}+C l^{2},\qquad
A=\frac{\lambda^{2}}{4a^{2}},\;B=\frac{\lambda^{2}}{4b^{2}},\;C=\frac{\lambda^{2}}{4c^{2}} .
\label{eq:ch04b-HL}
\end{equation}
这里的待定量是 $q=\sin^{2}\theta$，\textbf{不是}倒易量 $Q=1/d^{2}$，两者只差一个常数因子
$q=(\lambda^{2}/4)Q$：相减与相除时该因子都会约去，所以下面的差值关系对二者同样成立，但
$A,B,C$ 的数值必须与反算公式配套——本节取 $A=\lambda^{2}/(4a^{2})$，反算用
$a=\lambda/(2\sqrt{A})$、$b=\lambda/(2\sqrt{B})$、$c=\lambda/(2\sqrt{C})$（若改按
$Q=1/d^{2}$ 记账，系数就是 $1/a^{2}$、反算 $a=1/\sqrt{A}$，两套不能混用）。
于是“有关联但指标不同”的 $q$ 值之间存在严格的加和与倍数关系，例如
$q_{hkl}=q_{h00}+q_{0k0}+q_{00l}$，$q_{200}=4q_{100}$，$q_{300}=9q_{100}$，
$q_{hk0}=q_{h00}+q_{0k0}$ 等。\textbf{这些关系就是该方法的全部依据}：只要能从实测数据中识别出
哪些差值对应轴面衍射，三个参数就被确定（正交晶系要三条互不成比例的线才能定出 $A,B,C$
三个独立系数）。
（此处已对照原页面图像（书内 p.266 / PDF p.285）：原书即令 $q_{hkl}=\sin^{2}\theta_{hkl}$，
并把正交晶系的方程变换为 $q_{hkl}=\sin^{2}\theta=Ah^{2}+Bk^{2}+Cl^{2}$，其中
$A=\lambda^{2}/4a^{2}$、$B=\lambda^{2}/4b^{2}$、$C=\lambda^{2}/4c^{2}$，与本式一致；原书同页给出的
加和与倍数关系为 $q_{hkl}=q_{h00}+q_{0k0}+q_{00l}=q_{hk0}+q_{00l}=q_{h0l}+q_{0k0}=\cdots$ 及
$q_{1kl}-q_{0kl}=q_{100}$、$q_{2kl}-q_{0kl}=4q_{100}$、$q_{3kl}-q_{0kl}=9q_{100}$，
本式所写的 $q_{200}=4q_{100}$、$q_{300}=9q_{100}$ 是其在 $k=l=0$ 时的特例。原书未另行定义
倒易量 $Q=1/d^{2}$；“$q$ 与 $Q$ 相差 $\lambda^{2}/4$、两套系数不能混用”是本节的提醒。）
\end{principle}

\begin{measure}
处理过程：读出各峰衍射角，求 $\sin^{2}\theta$，由小到大排列后\textbf{依次相减}
$\Delta\sin^{2}\theta_{ij}=\sin^{2}\theta_{j}-\sin^{2}\theta_{i}$，统计差值出现的频度，并观察
差值之间的整数比关系。\textbf{选出现频度最高的差值作为轴面衍射的 $\sin^{2}\theta$}（即 $A$、
$B$ 或 $C$ 中的某一个），再用它去除其余差值：若商中出现 $1,4,9,\dots$ 即
$1:4:9$ 关系，说明该差值确实对应 $100$（或 $010$、$001$）及其高次衍射，据此把它定为
$q_{100}$；接着寻找可能的 $q_{010}$ 与 $q_{001}$，对全部线进行拟合，重复数次直到所有衍射线
都能被指标化。对于单斜和三斜晶系，处理过程更为烦琐，需要更多尝试，借助计算机程序是更有效的
途径。
\end{measure}

\begin{dataflow}
$\sin^{2}\theta$ 升序表 $\rightarrow$ 相邻（及非相邻）差值表 $\rightarrow$ 差值频度统计
$\rightarrow$ 高频差值作为轴面衍射候选 $\rightarrow$ 用 $1:4:9$ 关系验证候选 $\rightarrow$ 确定
$A$、$B$、$C$ $\rightarrow$ 反算 $a,b,c$ $\rightarrow$ 用全部线回代检验（所有线都应落在
允许误差内）。频度统计之所以有效，是因为轴面衍射的 $q$ 值会在很多差值组合中重复出现，
而非轴面反射的组合关系没有这种普遍性。
\end{dataflow}

\subsubsection{常见指标化程序的原理}

20 世纪 60 年代以来出现多种计算机指标化程序，主要计算原理有晶带法、面指数尝试法、连续二分法
和蒙特卡罗方法。\src{书内 p.267 / PDF p.286}

\paragraph{TREOR——晶面指数尝试法}1985 年 Werner 等人发表基于面指数尝试法的程序 TREOR，
其核心是在倒易空间中的半穷举搜索：选若干低角度衍射线作为基本晶面的衍射，赋予它们指定范围内
的小整数指标 $h,k,l$（例如 $|h|,|k|,|l|<3$），代入晶面间距与晶胞参数的关系式求解方程、算出
晶胞参数；再用所得参数尝试指标化其余的衍射线，并寻找最可能的解。以单斜晶系为例，关系式中含有
四个待求晶胞参数，故取四条相互独立的低角度衍射线进行尝试，算出倒易格子参数后逐一计算第 6、7
等衍射峰并与实验值比较，直到几乎所有衍射线都被纳入且对照合理。TREOR 对单斜及更高对称性晶系
成功率在 $95\%$ 以上，优点是计算速度快、可指定容忍误差与允许不能指标化的衍射数目；缺点是
对三斜晶系处理难度大，不能自动指认空间群与带心晶格，且品质因子 $M_{20}$ 与 $F_{N}$ 可能被低估。

\paragraph{DICVOL——二分法}1972 年 Louër 首次提出二分法指标化方法，据此发展出程序 DICVOL。
它从正空间出发，以晶胞边长及夹角为变量，在给定区间内用二分法逐步缩小晶胞参数范围。以立方晶系
为例：给定 $a$ 在 $a_{0}$ 与 $a_{M}$ 之间、初始搜索步长 $p$（如 $0.5$ Å），从 $a_{0}$ 算到
$a_{M}$，每一步计算
\begin{equation}
Q^{-}(hkl)=\frac{h^{2}+k^{2}+l^{2}}{\left[a_{0}+(n+1)p\right]^{2}},\qquad
Q^{+}(hkl)=\frac{h^{2}+k^{2}+l^{2}}{\left[a_{0}+np\right]^{2}} ,
\label{eq:ch04b-dicvol}
\end{equation}
并把每步结果与全部实验 $Q$ 值（一般取前 $20$ 个）比较，若在误差范围内满足
$Q^{-}\le Q_{\mathrm{obs}}\le Q^{+}$，就把该区间分成两半、步长减半（$0.25$ Å）继续计算与测试；
找到更接近的数值范围后再二分（$0.125$ Å）。连续做 6 次二等分，可将晶胞参数的变化范围限制在
$0.5/2^{6}=0.0078$ Å 之内。DICVOL 的优点是正空间穷举搜索原则上不会遗漏，且先处理体积较小的
晶胞，通常可避免晶胞体积加倍的假解；缺点是三斜与单斜晶系计算时间长。
（此处已对照原页面图像（书内 p.268 / PDF p.287）：原书即印
$Q_{-}(hkl)=(h^{2}+k^{2}+l^{2})/[a_{0}+(n+1)p]^{2}$ 与
$Q_{+}(hkl)=(h^{2}+k^{2}+l^{2})/[a_{0}+np]^{2}$，判别式印为
$Q_{-,hkl}\le Q_{\mathrm{obs}}\le Q_{+,hkl}$；式~\equref{eq:ch04b-dicvol} 与原书逐字一致，
$n$ 与 $n+1$ 的对应关系就是原书的印法。）

\subsubsection{指标化结果的判断：品质因子}

指标化在数学上是多解问题，通常可以得到多种结果。判断哪一个最可能正确，最常用的两个品质因子
是 de Wolff 提出的 $M_{20}$ 与 Smith 提出的 $F_{N}$。\src{书内 p.268—269 / PDF p.287—288}

\begin{principle}
$M_{20}$ 由 “第 20 条线的 $Q$ 值” 与 “平均偏差、计算线条数” 构成：
\begin{equation}
M_{20}=\frac{Q_{20}}{2\,\overline{\left|\Delta Q\right|}\,N_{20}} ,
\label{eq:ch04b-M20}
\end{equation}
式中 $Q_{20}$ 为所观察到的（且已指标化的）第 20 条衍射线的 $Q$ 值；$N_{20}$ 为计算 $Q$ 值到
$Q_{20}$ 为止所得出的不同 $Q$ 值的个数；$\overline{\left|\Delta Q\right|}$ 为 $Q$ 的观察值与
计算值之间的平均偏差（指指标化后各线偏差绝对值的平均）。观察线条数与计算线条数的比值反映了
“计算出的线条中有多少能在实验中被看到”。$M_{20}$ 越大越好。\\
$F_{N}$ 定义为两项的乘积，前一项与准确性有关、后一项与数据完备性有关：
\begin{equation}
F_{N}=\frac{1}{\overline{\left|\Delta2\theta\right|}}\cdot\frac{N_{\mathrm{obs}}}{N_{\mathrm{poss}}} ,
\label{eq:ch04b-FN}
\end{equation}
式中 $N_{\mathrm{obs}}$ 为观察到的衍射线数目，$N_{\mathrm{poss}}$ 为到第 $N$ 条观察衍射线为止
按所得晶胞与消光规律推算出的\textbf{可能}衍射线数目，$\overline{\left|\Delta2\theta\right|}$
为计算值与观察值之间偏差绝对值的平均（单位与 $2\theta$ 一致，通常取度）。Smith 建议 $N$ 取
$30$；若衍射线不够 $30$ 条，则取全部衍射线数目。$F_{N}$ 值越大，结果越可靠。
\textbf{比值的方向不能写反}：$N_{\mathrm{obs}}\le N_{\mathrm{poss}}$，完备性因子不大于 $1$；
同一个晶胞若额外预言出一批实验上并未出现的线，$N_{\mathrm{poss}}$ 增大而 $F_{N}$ 必须
\textbf{下降}。若把比值写成 $N_{\mathrm{poss}}/N_{\mathrm{obs}}$，就变成奖励「多预言、不被
观测」的晶胞，正好与指标化的目标相反。
此处已对照原页面图像（书内 p.268—269 / PDF p.287—288）：原书印
$M_{20}=\dfrac{Q_{20}}{2\bar{\varepsilon}N_{20}}$，其中
$\bar{\varepsilon}=\dfrac{1}{20}\sum_{i=1}^{20}|\Delta Q_{i}|$ 为 $Q$ 的观察值与计算值之间的
平均偏差，与上式 $\overline{|\Delta Q|}$ 同义，分母确为 $2\bar{\varepsilon}N_{20}$；
$N_{20}$ 原书措辞为「计算 $Q$ 值到 $Q_{20}$ 时所得出的不同 $Q$ 值的个数」，与本节一致。
$F_{N}$ 原书印 $F_{N}=\dfrac{1}{|\overline{\Delta2\theta}|}\cdot\dfrac{N}{N_{\text{预测}}}$，
并明确「$N_{\text{预测}}$ 为到第 $N$ 条观察衍射线的可能衍射线的数目」——即
$N_{\mathrm{poss}}$ 取「到第 $N$ 条观察线为止」，本节写法（比值取
$N_{\mathrm{obs}}/N_{\mathrm{poss}}$）与原书一致；原书另注 Smith 建议 $N$ 取 $30$，
不足 $30$ 则取全部衍射线数目。
\end{principle}

\begin{table}[tbp]
  \centering
  \caption{品质因子的判据。两个因子分别从“指标化精度”和“数据完备性”两个角度衡量同一套解，
    单独使用任何一个都可能误判，通常要求两者同时通过。\src{据 S1（上册）§4.6.4，书内 p.268—269 / PDF p.287—288 归纳。}}
  \label{tab:ch04b-FOM}
  \begin{tabular}{@{}p{16mm}p{62mm}p{46mm}@{}}
    \toprule
    因子 & 物理含义 & 判据 \\
    \midrule
    $M_{20}$ & 第 20 条线的 $Q$ 值与平均偏差、计算线条数之比；综合反映偏差大小与“计算线是否
    都能被观察到” & 不能指标化的线条少于 $2$ 条且 $M_{20}>10$ 时结果基本正确；
    $M_{20}<6$ 时值得怀疑；$M_{20}<3$ 时几乎没有意义 \\
    $F_{N}$ & 平均 $2\theta$ 偏差的倒数（准确性）与“观察线数／可能线数”（完备性）的乘积 &
    越大越可靠；多数正确解 $F_{N}>10$，但很难说某个具体数值就一定对应正确解 \\
    \bottomrule
  \end{tabular}
\end{table}

\begin{pitfall}
\textbf{指标化不唯一：品质因子必须与化学约束共同判定。}当两组解的 $M_{20}$ 与 $F_{N}$ 相近时，
应优先选择：（1）有更高对称性的晶胞；（2）体积最小的晶胞；（3）理论衍射线数目少的晶胞。
除此之外还必须综合判断以下六点——其中任何一点不过关，都不能仅凭品质因子宣布结果正确：
（1）最初输入的全部衍射线是否都能被指标化，未被指标化的线要找出解释（是弱峰、杂相、还是
K$\beta$／钨污染线）；（2）未参与最初指标化的低角度弱峰和高角度峰能否基于所得晶胞给以恰当的
指标，若可以则结果更可靠；（3）指标化结果中是否含有明确的系统消光规律（有则更可靠，见表
~\tabref{tab:ch04b-extinction}）；（4）检索数据库能否找到具有类似结构的化合物，以及是否与
电子衍射结果一致；（5）若已知分子式和样品密度，可根据单胞内化学式单位数 $Z$ 是否为整数来判断
——对有机分子，晶胞体积 $V$ 与晶胞内分子数 $Z$ 之间有近似关系 $V\approx18Z N_{\mathrm{NHA}}$
（$N_{\mathrm{NHA}}$ 为分子中非氢原子 C、N、O 等的总数），这条判据对无机盐与氧化物需改用密度
与式量的关系；（6）用多种计算程序是否都得到相同结果。最后还要检验所得的键长、键角在化学上
是否合理，并在指标化结果基础上做全谱拟合，看最终能否解出合理的晶体结构。
此处已对照原页面图像（书内 p.269 / PDF p.288）：原书印 $V\approx18ZN_{\mathrm{NHA}}$，
$N_{\mathrm{NHA}}$ 为分子中非氢原子（C、N、O 等）的总数——系数 $18$ 与下标写法均照原书
（原书下标为 $\mathrm{NHA}$，本稿此前写作 $\mathrm{HA}$，已改齐）；原书同页另印
$M_{20}\ge10$ 结果基本正确、$M_{20}<6$ 值得怀疑、$M_{20}<3$ 几乎没有意义，以及
$F_{N}$ 多数正确解在 $10$ 以上，并给出「优先选择更高对称性、体积最小、理论衍射线数目小」
三条取舍顺序。
\end{pitfall}

\begin{example}[读谱实例：指标化多解与化学约束]
\begin{obs}
一组示意数据（构造，用于说明判据的使用）：某未知粉晶，化学分析给出
$n(\mathrm{Mg}):n(\mathrm{Al})\approx1:2$，$20$ 条衍射线。用 TREOR 得到两组候选解：
解 A 为立方，$a=8.08$ Å，$M_{20}=22$，全部 $20$ 条线均可指标化，且消光规律与面心格子一致；
解 B 为立方，$a=4.04$ Å（体积为 A 的 $1/8$），$M_{20}=19$，但有 $4$ 条低角弱峰无法指标化。
\end{obs}
\begin{infer}
两组解的 $M_{20}$ 都超过 $10$，\textbf{单凭品质因子无法取舍}。再引入化学与晶体学约束：
按密度与式量估算单胞内化学式单位数，解 A 给出整数 $Z=8$，解 B 给出 $Z=1$，且 $Z=1$ 的晶胞
容纳不下该组成的阳离子配位环境；同时解 A 能指标化全部弱峰并给出明确的系统消光规律，解 B
不能。因此取解 A，解 B 是解 A 的“体积亚倍”假解——这正是指标化多解性中最常见的一类。
\end{infer}
\textbf{判据小结}：品质因子回答“这套指标在数值上是否自洽”，化学约束回答“这个晶胞在物理上
是否可能存在”，两者缺一不可。
\end{example}

\begin{crosscheck}
指标化结果的外部校验途径：（1）与电子衍射（选区电子衍射给出三维倒易信息，可验证点阵类型与
晶胞尺寸）；（2）与已知同构化合物的 PDF 卡片、单晶结构数据对照，看晶胞参数与消光规律是否
一致；（3）与化学分析、密度测量给出的 $Z$ 相互印证；（4）用不同原理的程序（尝试法与二分法）
互算；（5）最终的检验是把指标化所得晶胞参数、样品组成与已知原子参数正算出衍射图，与实验
衍射图逐峰对照——能对上才说明指标化与结构模型都成立。
\end{crosscheck}

\begin{methodsum}
指标化是结构分析的第一步：它同时给出每个峰的晶面指标、晶胞参数 $(a,b,c,\alpha,\beta,\gamma)$、
晶系、点阵类型与可能的空间群；用一套晶胞参数指标化全部衍射线是确认单相的重要依据；它也是
晶胞参数精修与全谱拟合（Rietveld）的前提。立方晶系可用解析法（$Q$ 比值整数化 + 消光判据）
手工完成；四方、六方、正交晶系用 Hesse--Lipson 差值法；低对称晶系必须借助 TREOR、DICVOL、
ITO、McMaille 等程序。无论用哪种方法，结果都必须通过品质因子（$M_{20}$、$F_{N}$）与化学约束
（$Z$ 为整数、密度一致、键长键角合理、消光规律明确）的双重检验才能采用。
\end{methodsum}
